% !TeX root = ../main.tex

\chapter{Introduction}
\label{chap:introduction}

\section{Project Overview}

The Multi-Rail Isolated Power Supply Monitoring System is a hardware and
software system developed to provide centralized monitoring of multiple
electrically isolated power-supply rails. The system measures both output
voltage and load current, provides configurable alarm
limits, displays operating status and measurement data via a local display or ethernet for integration with an EPICS-based control system.

The completed system monitors six independent power-supply rails:

\begin{itemize}
    \item +5VA
    \item +5VB
    \item +5V
    \item -5V
    \item +15V
    \item -15V
\end{itemize}

Because the monitored supplies may be electrically isolated from one another,
the monitoring system must preserve that isolation while still allowing all six
channels to communicate with a common controller. Each supply rail therefore
uses an independent voltage/current monitoring channel and an isolated digital
communication path.

The primary controller is an Arduino GIGA R1 WiFi. The controller acquires
measurement data from the individual monitor channels, converts the data to
engineering units, applies the appropriate rail polarity, evaluates configured
alarm limits, and presents the resulting information on a GIGA Display Shield.
An Ethernet interface provides network communication between the embedded
controller and an EPICS Input/Output Controller (IOC). The EPICS IOC then makes
the measurements, alarms, and system status available to a CS-Studio Phoebus
operator interface.


\section{Project Motivation}

Power-supply systems containing multiple isolated voltage rails present a
monitoring challenge because each circuit associated with one supply
can't share a direct electrical reference with another supply.
Connecting conventional measurement or communication circuits across these
references can break isolation between the rails.

A centralized monitor is desirable because it allows the user to observe all supply voltages and currents from a single interface,
identify abnormal operating conditions, detect communication or sensor faults,
and integrate the power-supply status into a larger control system.

The design developed for this project addresses these requirements by
performing voltage and high-side current measurement locally on each monitored
rail and transferring the resulting information across electrically isolated
digital interfaces. The controller receives measurement data from all
six rails without requiring their power references to be directly connected
together.


\section{System Architecture}

The monitoring system is divided into three primary functional areas:
the power-supply monitoring hardware, the microcontroller, and the
supervisory control interface.

The power supply monitoring hardware contains the current-sense resistors,
LTC2945 PMICs, and digital isolation circuitry. Each LTC2945 measures the supply
voltage and the differential voltage developed across a high-side current-sense
resistor. The resulting digital measurement information is sent over the I\textsuperscript{2}C bus while maintaining the required electrical isolation.

The Arduino GIGA R1 WiFi provides the mcu processing functions. Its
firmware initializes the monitoring channels, acquires voltage and current
measurements, detects communication errors, converts raw measurements to
engineering units, applies channel-specific polarity and scaling, evaluates
alarm limits, and updates the local display. Configuration information,
including network parameters and alarm limits, is retained in persistent
storage so that required settings can survive a controller reset or power
cycle.

Network communication is provided through an Ethernet interface using the PSC
protocol. The EPICS IOC communicates with the embedded controller using
PSCDriver and converts the transmitted values into EPICS process variables.
CS-Studio Phoebus provides the remote operator interface for viewing
measurements, alarm conditions, connection status, and system state.


\section{Design Objectives}

The major design objectives for the Multi-Rail Isolated Power Supply Monitoring
System are to:

\begin{enumerate}
    \item Measure the voltage of six independent power-supply rails.

    \item Measure the load current of each rail using high-side current sensing.

    \item Preserve electrical isolation between independently referenced
    power-supply rails and the common controller.

    \item Convert acquired data into meaningful engineering units and correctly
    represent both positive and negative voltage rails.

    \item Detect I\textsuperscript{2}C communication failures and prevent invalid
    sensor data from being presented as valid measurements.

    \item Compare voltage and current measurements against configurable upper and
    lower alarm limits.

    \item Provide a local graphical display of voltage, current, alarm status,
    communication status, and network state.

    \item Retain configuration and alarm settings through controller resets and
    power cycles.

    \item Provide Ethernet communication for remote monitoring and configuration.

    \item Integrate the monitoring system with an EPICS IOC and CS-Studio Phoebus
    operator interface.

    \item Maintain the hardware, firmware, and control-system software under
    revision control so that changes made during design, implementation, and
    testing remain traceable.
\end{enumerate}


\section{Project Scope}

The scope of the project includes the design and implementation of the sensor
interface hardware, Arduino GIGA controller hardware, embedded firmware,
Ethernet communication interface, EPICS IOC, and operator-interface software.
The hardware portion includes the six voltage/current monitoring channels,
current-sense resistors, digital isolation circuitry, controller interface,
Ethernet circuitry, and supporting power-supply circuitry.

The software portion includes the embedded acquisition and alarm-processing
firmware, persistent configuration storage, local display interface, PSC
network communication, EPICS IOC configuration and database records, and the
CS-Studio Phoebus operator interface.

Design files and software associated with the project are maintained in the
project Git repository:

\begin{center}
    \url{https://github.com/capotosto/PS_Monitor}
\end{center}

The repository provides revision control for the hardware design, firmware,
EPICS IOC files, documentation, and other project artifacts generated during
development.


\section{Report Organization}

This report documents the development of the monitoring system from preliminary
design through implementation and testing. The preliminary-design chapter
defines the major functional blocks, presents the system architecture,
documents supporting engineering calculations, includes the hardware
schematics and bills of materials, and describes the primary firmware and
control-system software flows.

Subsequent chapters document prototyping and development activities, the EPICS
IOC implementation, and the CS-Studio Phoebus operator interface. Supporting
project-management and design documents developed during the earlier phases of
the capstone project are retained as part of the overall project documentation
and provide traceability between the original requirements, design decisions,
implementation, and final system.